\documentclass[12pt]{article}
\usepackage[utf8]{inputenc}
\usepackage[T1]{fontenc}
\usepackage[margin=1in]{geometry}
\usepackage{setspace}
\usepackage{xcolor}
\usepackage[colorlinks=true, urlcolor=blue!60!black, linkcolor=blue!60!black]{hyperref}
\setlength{\parindent}{2em}
\setlength{\parskip}{1.25em}
\renewcommand{\baselinestretch}{1.0}
\begin{document}
\begin{center}
{\LARGE \textit{BMW of North America, Inc. v. Gore}, 517 U.S. 559 (1996)\par}
\vspace{0.8em}
{\large Case Brief --- Four-Page Edition\par}
\end{center}
\vspace{1em}

\section{Memory Jogger}

For the first time in modern jurisprudence, the Court held a state punitive damages award---\$2 million on \$4,000 in actual damages for failing to disclose a \$601.37 presale paint repair---``grossly excessive'' under the Fourteenth Amendment's Due Process Clause, establishing three guideposts (reprehensibility, ratio, and comparable sanctions) for reviewing punitive awards (\href{https://www.courtlistener.com/opinion/118026/bmw-of-north-america-inc-v-gore/}{517 U.S. 559}).

\section{Detailed Case Facts}

In 1990 Dr. Ira Gore, Jr., bought a new BMW sedan for \$40,750.88 in Birmingham, Alabama. Nine months later a detailing shop discovered that the car had been repainted at BMW's preparation center, presumably after acid-rain damage in transit. Since 1983 BMW had followed a nationwide policy: cars whose repair costs exceeded 3\% of MSRP were sold as used; cars below that threshold were repaired and sold as new without disclosure. Gore's repainting cost about 1.5\% of MSRP.

Gore sued for fraudulent suppression under Alabama law. A former dealer testified that repainting cut the car's value by about 10\%, or \$4,000. Gore showed that BMW had sold 983 refinished cars nationwide as new, 14 of them in Alabama, and asked for \$4 million in punitive damages (\$4,000 times about 1,000 cars). BMW answered that its 3\% threshold complied with the disclosure statutes of roughly 25 states and that it had never before been held liable for punitive damages for the policy. Five days after the verdict, BMW adopted full disclosure nationwide. The case arrived amid the 1990s tort-reform debate; \textit{Haslip} and \textit{TXO} had recognized a due process limit on ``grossly excessive'' awards in theory but upheld both awards.

\section{Procedural History}

An Alabama jury awarded \$4,000 in compensatory and \$4 million in punitive damages; the trial judge denied BMW's motion to set aside the punitive award. The Alabama Supreme Court held that the jury had improperly multiplied by out-of-state sales, disclaimed reliance on them, and remitted the award to \$2 million. The U.S. Supreme Court granted certiorari and, on May 20, 1996, reversed and remanded.

\section{Judicial Votes}

Reversed and remanded, \textbf{5--4}, across conventional ideological lines.
\begin{itemize}
\item \textbf{Majority:} Justice Stevens, joined by Justices O'Connor, Kennedy, Souter, and Breyer.
\item \textbf{Concurrence:} Justice Breyer, joined by Justices O'Connor and Souter.
\item \textbf{Dissent:} Justice Scalia, joined by Justice Thomas.
\item \textbf{Dissent:} Justice Ginsburg, joined by Chief Justice Rehnquist.
\end{itemize}

\section{Holding}

The \$2 million award was grossly excessive and violated the Due Process Clause. A State may not use punitive damages to punish a defendant's lawful conduct in other jurisdictions, and a defendant must receive fair notice of both the conduct subject to punishment and the severity of the penalty. Courts assess excessiveness by (1) the degree of reprehensibility, (2) the ratio between the actual or potential harm and the award, and (3) the difference between the award and civil or criminal penalties for comparable misconduct. The extended extraterritoriality discussion, footnote 21's admission of out-of-state conduct as evidence of reprehensibility, and the remarks about higher ratios for small or hard-to-detect injuries are arguably dicta.

\section{Analysis of Opinions}

\textbf{Majority (Stevens).} Because state disclosure laws vary and reasonable lawmakers disagree, Alabama may protect its own consumers but may not punish conduct lawful elsewhere or force a nationwide policy change; the award had to be judged against BMW's 14 Alabama sales. Reprehensibility, ``[p]erhaps the most important indicium,'' was low: the harm was purely economic, safety was never threatened, BMW could reasonably read the state statutes as safe harbors, and there was no deliberate falsehood. The 500:1 ratio (35 times even the \$56,000 statewide harm) dwarfed \textit{Haslip} and \textit{TXO}. And Alabama's maximum civil penalty for deceptive practices was \$2,000, so nothing gave BMW fair notice of a multimillion-dollar sanction.

\textbf{Concurrence (Breyer).} The strong presumption of validity for fairly produced state judgments is overcome because Alabama's regime lacked constraining standards: its statute was too broad, the \textit{Green Oil} factors were applied without force, no economic theory explained the figure, history offered no precedent, and no legislative cap applied.

\textbf{Dissent (Scalia).} Due process guarantees only procedure, not reasonable awards. The extraterritoriality discussion is ``the purest dicta,'' and the guideposts ``mark a road to nowhere,'' federalizing state tort law without workable rules.

\textbf{Dissent (Ginsburg).} The Alabama Supreme Court had already cured the jury's nationwide multiplier and reviewed the award under factors approved in \textit{Haslip}. With no bright-line test, the Court will stand alone policing state awards; state legislatures, already enacting caps and other reforms, should handle the problem.

\section{Examples: Future Applications}

\begin{enumerate}
\item \textbf{Same Side (Excessive): Insurer Underpayment.} A policyholder with \$10,000 in economic loss wins \$15 million (1,500:1) based on nationwide claims volume, while the state fine is \$5,000 per violation. Purely economic harm, an extreme ratio, and tiny comparable fines make the award unconstitutional.
\item \textbf{Same Side (Excessive): Battery-Draining Software.} A buyer recovers \$200 in compensatory and \$500,000 in punitive damages (2,500:1) for an undisclosed telemetry feature, with civil penalties capped at \$2,500. All three guideposts point to excessiveness.
\item \textbf{Opposite Side (Permissible): Concealed Drug Danger.} A drugmaker hides data showing fatal liver failure and falsifies FDA submissions; a \$4 million punitive award on \$1 million in compensatory damages (4:1) survives given extreme reprehensibility and severe criminal sanctions.
\item \textbf{Opposite Side (Permissible): Winter Utility Shutoff.} A landlord cuts heat and water to force out an elderly tenant; \$35,000 in punitive damages on \$5,000 in losses (7:1) aligns with \$25,000 in authorized fines and passes.
\item \textbf{Fence-Sitter: Sabotage With Small Actual Harm.} A competitor dumps debris to sabotage a bid, causing \$500 in cleanup costs but threatening a \$200,000 contract. A \$150,000 award is 300:1 to actual harm but under 1:1 to potential harm, and \textit{Gore} allows higher ratios for intentional misconduct causing small actual damages.
\end{enumerate}

\section{Critique}

\textbf{Originalist view.} Echoing Scalia, critics argue that \textit{Gore} revives \textit{Lochner}-style economic substantive due process without textual or historical mooring, and that the Court, in the name of federalism, federalized state tort law and displaced state courts and legislatures already enacting reforms.

\textbf{Progressive view.} Tying punitive awards to single-plaintiff ratios undermines deterrence of widespread, low-dollar fraud that few victims detect, and ``fair notice'' rings hollow for sophisticated corporations. The unstable guideposts produced unpredictable outcomes until \textit{State Farm v. Campbell} tightened them into a presumptive single-digit ceiling.

\section{Key Quotations}

\begin{quote}
``Elementary notions of fairness enshrined in our constitutional jurisprudence dictate that a person receive fair notice not only of the conduct that will subject him to punishment, but also of the severity of the penalty that a State may impose.''
\end{quote}
\noindent --- Justice Stevens, Opinion of the Court, 517 U.S. at 574.

\begin{quote}
``Punitive damages may properly be imposed to further a State's legitimate interests in punishing unlawful conduct and deterring its repetition. ... Only when an award can fairly be categorized as `grossly excessive' in relation to these interests does it enter the zone of arbitrariness that violates the Due Process Clause of the Fourteenth Amendment.''
\end{quote}
\noindent --- Justice Stevens, Opinion of the Court, 517 U.S. at 568.

\begin{quote}
``In most cases, the ratio will be within a constitutionally acceptable range, and remittitur will not be justified on this basis. When the ratio is a breathtaking 500 to 1, however, the award must surely `raise a suspicious judicial eyebrow.'\,''
\end{quote}
\noindent --- Justice Stevens, Opinion of the Court, 517 U.S. at 583.

\end{document}
